% =============================================================
%  6. CONCLUSIONES
% =============================================================
\renewcommand{\seccorta}{Conclusiones}
\section{Conclusiones}

\begin{frame}{Conclusiones}
  \begin{itemize}
    \item El enlace de 70 m alcanzó la máxima modulación (256QAM, $\approx$450 Mbps de
          capacidad) con solo $-4$ dBm de potencia.
    \item El \textbf{multipath} del pasillo generó desbalance entre polarizaciones y jitter;
          la \textbf{obstrucción de Fresnel} por personas produjo caídas de hasta 28 dB.
    \item La \textbf{modulación adaptativa} mantiene el enlace ante desalineación u
          obstrucciones, sacrificando capacidad.
    \item El \textbf{ancho de canal} es un compromiso: 60 MHz maximiza capacidad; 10 MHz
          reduce ruido y jitter.
    \item Las herramientas embebidas tienen límites (Speed Test $\approx$267 Mbps):
          para medir el enlace real hace falta iPerf3 con PCs en ambos extremos.
    \item \completar[conclusiones del otro grupo]
  \end{itemize}
\end{frame}

\begin{frame}[plain]
  \centering\vfill
  {\usebeamercolor[fg]{title}\Huge\bfseries ¡Gracias!}\\[8pt]
  {\color{naranja}\large ¿Preguntas?}
  \vfill
\end{frame}
